\documentclass[11pt]{amsart}
\usepackage[a4paper,margin=1in]{geometry}
\usepackage{amsmath,amssymb,amsthm,mathtools}
\usepackage{booktabs,longtable,array}
\usepackage{microtype}
\usepackage{xcolor}
\usepackage{hyperref}
\hypersetup{colorlinks=true,linkcolor=blue!45!black,citecolor=blue!45!black,
  urlcolor=blue!45!black,
  pdftitle={Uniformity, Prime Support, and Reachable Lattices in Approaches to the abc Conjecture},
  pdfauthor={ChatGPT}}
\newtheorem{theorem}{Theorem}[section]
\newtheorem{proposition}[theorem]{Proposition}
\newtheorem{lemma}[theorem]{Lemma}
\newtheorem{corollary}[theorem]{Corollary}
\theoremstyle{definition}
\newtheorem{definition}[theorem]{Definition}
\theoremstyle{remark}
\newtheorem{remark}[theorem]{Remark}
\newcommand{\N}{\mathbb N}
\newcommand{\Q}{\mathbb Q}
\newcommand{\R}{\mathbb R}
\newcommand{\Z}{\mathbb Z}
\newcommand{\OO}{\mathcal O}
\newcommand{\rad}{\operatorname{rad}}
\newcommand{\Span}{\operatorname{span}}
\newcommand{\GL}{\operatorname{GL}}
\newcommand{\Mat}{\operatorname{Mat}}
\newcommand{\den}{\operatorname{den}}
\newcommand{\num}{\operatorname{num}}
\newcommand{\eps}{\varepsilon}
\newcommand{\loglog}{\log\log}
\newcommand{\height}{\mathrm h}
\newcommand{\pospart}[1]{\left(#1\right)_+}

\title[Uniformity and prime support in $abc$ research]{Uniformity, Prime Support, and Reachable Lattices\\
in Approaches to the $abc$ Conjecture}
\author{ChatGPT}
\date{August 30, 2026}
\subjclass[2020]{11D75, 11N25, 11G05, 11J86, 68V15}
\keywords{abc conjecture, smooth numbers, prime support, effective integral points,
local lattices, formal verification}

\begin{document}
\begin{abstract}
We examine precise uniformity requirements in several approaches to the $abc$
conjecture. An elementary encoding bounds the number of integers $n\le N$
with all prime divisors at most $Y$ and at most $w$ distinct prime divisors by
$(1+Y\lfloor\log_2N\rfloor)^w$. Combined with an unconditional
short-interval theorem of Younis, this shows that proposed smooth-number
moment estimates in a recent claimed disproof are false in their stated
parameter range. We also give the exact coupled prime-exponent formulation
of the $abc$ height defect and an unbounded dyadic obstruction to estimating
its endpoints separately. For local tensor constructions, we prove a finite
determinant certificate for the Haar measure of the lattice generated by
reachable vectors, and distinguish independent actions from simultaneous
ones. Effective integral-point bounds constrain a specified Pell--Chebyshev
residual family, while an explicit lower bound for the denominator of the
Frey $j$-invariant limits the applicability of recent small-denominator
results. Elementary arithmetic, counting, and linear-algebra components are
checked in Lean; external analytic inputs and remaining interfaces are
identified separately. Neither a proof nor a disproof of $abc$ is claimed.
\end{abstract}
\maketitle

\section{The target and the role of the results}

Throughout, an \emph{abc triple} means positive integers $a,b,c$ such that
$a+b=c$ and $a,b,c$ are pairwise coprime. Write
\[
 r=\rad(abc),\qquad R=\log r,\qquad h=\log c.
\]
Here $\rad(n)$ is the product of the distinct prime divisors of $n$.
Pairwise coprimality is equivalent to $\gcd(a,b,c)=1$ under the additive
relation. The target is the standard assertion
\begin{equation}\label{eq:abc}
 \forall\eps>0\quad\exists C_\eps\in\R\quad
 \forall(a,b,c):\qquad h\le(1+\eps)R+C_\eps.
\end{equation}
Exponentiation gives the usual multiplicative constant. The repository's
\texttt{ABCConjecture} uses $\log\max(a,b,c)$, which equals $h$ here.
Its definition is not changed in this work.

The quantifier order in \eqref{eq:abc} is essential. Its negation requires
one fixed positive $\eps$ and violations for every additive constant, hence
at unbounded height. A finite set of high-quality triples cannot supply that
negation: it is absorbed into $C_\eps$. Likewise, a sparse exceptional-set
bound does not imply that the exceptional set is finite.

Our results concern four specific interfaces. Section~\ref{sec:smooth}
uses finite prime support to test a proposed analytic construction.
Section~\ref{sec:signed} records the cancellation that an endpoint approach
must preserve. Section~\ref{sec:lattice} supplies a local generator
certificate relevant to possible-image constructions. Sections~\ref{sec:pell}
and \ref{sec:frey} examine effective arithmetic geometry and recent
literature. These components are useful tests and partial theorems; none
constructs the missing uniform bound in \eqref{eq:abc}.

The new statements below were proved mathematically before their
corresponding Lean implementations. Independent cross-reviews checked the
arguments, source hypotheses, and formal scope. Classical published
theorems may be used in the paper, but are not silently represented as
proved Lean declarations. The formal scope is specified in
Section~\ref{sec:formal}. The literature cutoff is August 30, 2026.

\section{Prime support in short intervals}\label{sec:smooth}

Let $\omega(n)$ be the number of distinct prime divisors of $n$, and let
$P^+(n)$ be its greatest prime divisor, with $P^+(1)=1$.
For integers $N,Y\ge1$ and $w\ge0$, define
\[
 \mathcal A(N,Y,w)=
 \{1\le n\le N:P^+(n)\le Y,\ \omega(n)\le w\}.
\]

\begin{theorem}[Finite prime-power encoding]\label{thm:encoding}
With $L=\lfloor\log_2N\rfloor$, one has
\begin{equation}\label{eq:encoding}
 |\mathcal A(N,Y,w)|\le(1+YL)^w.
\end{equation}
\end{theorem}
\begin{proof}
Use the alphabet
\[
 \mathcal Q=\{1\}\cup\{q^e:1\le q\le Y,\ 1\le e\le L\}.
\]
It has at most $1+YL$ members; the bases need not be prime.
For $n\in\mathcal A(N,Y,w)$, write its actual prime factorization as
$n=\prod_{j=1}^t p_j^{e_j}$, with $p_1<\cdots<p_t$ and $t\le w$.
Since $2^{e_j}\le p_j^{e_j}\le n\le N$, each $e_j\le L$.
Thus every factor belongs to $\mathcal Q$. Append $w-t$ copies of $1$ to
obtain a word of length $w$. Its product recovers $n$, so different integers
have different encoded words. Counting all words proves \eqref{eq:encoding}.
The empty factorization covers $n=1$ and $w=0$; $N=1$ is also included.
\end{proof}

\begin{corollary}[A finite first-moment bound]\label{cor:moment}
Let $S$ be a finite set of distinct positive $Y$-smooth integers at most
$N$, and put $B=(1+Y\lfloor\log_2N\rfloor)^w$. Then
\begin{align}
 |\{n\in S:\omega(n)\le w\}|&\le B,\label{eq:lowcount}\\
 \sum_{n\in S}\omega(n)&\ge(w+1)\max\{|S|-B,0\}.
 \label{eq:momentfinite}
\end{align}
\end{corollary}
\begin{proof}
The first set is contained in $\mathcal A(N,Y,w)$. Its complement in $S$
has at least $\max\{|S|-B,0\}$ members, each with at least $w+1$ distinct
prime factors. Summation proves the second assertion.
\end{proof}

Suppose now that $y=y(x)\ge2$ satisfies
\begin{equation}\label{eq:yscale}
 \loglog x=o(\log y),\qquad \log y=o(\log x),
 \qquad u=\frac{\log x}{\log y}.
\end{equation}
For every fixed $\kappa>0$, Theorem~\ref{thm:encoding}, with
$N=\lfloor2x\rfloor$, $Y=\lfloor y\rfloor$ and
$w=\lfloor\kappa u\rfloor$, gives
\begin{equation}\label{eq:global-low}
 \#\{n\le2x:P^+(n)\le y,\ \omega(n)\le\kappa u\}
 \le x^{\kappa+o(1)}.
\end{equation}
Indeed,
\[
 w\log(1+YL)
 \le\kappa\frac{\log x}{\log y}
       \bigl(\log y+O(\loglog x)\bigr)
 =\kappa\log x+o(\log x).
\]
More generally the upper bound is $x^{o(1)}$ whenever
$w(\log y+\loglog x)=o(\log x)$.

\subsection{The external analytic input}
Write $\Psi(x,y)=\#\{1\le n\le x:P^+(n)\le y\}$.
Younis's unconditional Theorem~1.1 \cite{Younis}, p.~2, states that for
each fixed $17/30<\theta\le1$ there is $C_Y(\theta)>0$ such that
\begin{equation}\label{eq:younis}
 \frac{\Psi(x+h,y)-\Psi(x,y)}h
 =\frac{\Psi(x,y)}x
  \left(1+O_\theta\left(\frac{\log(u+1)}{\log y}\right)\right)
\end{equation}
uniformly for
\[
 x^\theta\le h\le x,\qquad
 \exp\!\left(C_Y(\theta)(\log x)^{2/3}(\loglog x)^{4/3}\right)
 \le y\le2x.
\]
Fix
\begin{equation}\label{eq:parameters}
 \frac{17}{30}<\theta<1,\quad C\ge C_Y(\theta),\quad
 h=x^\theta,\quad
 y=\exp\!\left(C(\log x)^{2/3}(\loglog x)^{4/3}\right).
\end{equation}
The lower bound on $C$ is part of the hypotheses, not a disposable constant.
Put
\[
 S_x=\{n\in(x,x+h]:P^+(n)\le y\},\qquad M_x=|S_x|.
\]
The classical long-interval estimate recalled in
\cite[equations (1.8)--(1.9), p.~415]{HT} yields
$\Psi(x,y)=x\rho(u)(1+o(1))$ here. Its range
$u\le(\log y)^{3/5-\eps}$ is satisfied, for example with $\eps=1/20$:
the power of $\log x$ on the right is $11/30$, larger than the power
$1/3$ in $u$.
Also $-\log\rho(u)=u\log u(1+o(1))=o(\log x)$.
Consequently \eqref{eq:younis} gives
\begin{equation}\label{eq:mass}
 M_x=h\rho(u)(1+o(1))=x^{\theta-o(1)}.
\end{equation}
Only this mass estimate is needed below; no short-interval moment formula
is assumed.

\begin{theorem}[A lower tail and a first-moment obstruction]\label{thm:tail}
For \eqref{eq:parameters} and every fixed $0<\kappa<\theta$,
\begin{equation}\label{eq:tail}
 \frac{\#\{n\in S_x:\omega(n)\le\kappa u\}}{M_x}
 \le x^{\kappa-\theta+o(1)}\longrightarrow0.
\end{equation}
In particular,
\begin{equation}\label{eq:mean}
 \liminf_{x\to\infty}\frac1{uM_x}
          \sum_{n\in S_x}\omega(n)\ge\theta.
\end{equation}
For every fixed $A>0$, the proportion of $n\in S_x$ with
$\omega(n)\le(\loglog x)^A$ tends to zero.
\end{theorem}
\begin{proof}
Every member of $S_x$ is at most $2x$. Divide \eqref{eq:global-low} by
\eqref{eq:mass} to obtain \eqref{eq:tail}. If its numerator is $B_x$, then
\[
 \sum_{n\in S_x}\omega(n)\ge\kappa u(M_x-B_x).
\]
The lower limit after division by $uM_x$ is at least each fixed
$\kappa<\theta$, and hence at least $\theta$. This argument does not
require a variable $\kappa$ or a uniform error term in that variable.
Finally,
\[
 (\loglog x)^A(\log y+\loglog x)=o(\log x).
\]
The last observation after \eqref{eq:global-low} bounds the corresponding
population by $x^{o(1)}$. Divide again by \eqref{eq:mass}.
\end{proof}

\begin{corollary}[A false intermediate assertion in a claimed disproof]
\label{cor:carella}
Lemma~4.2 of Carella \cite{Carella}, p.~8, equation~(4.5), is false in its
stated range. The same holds for the second-moment assertion in
Lemma~4.4, p.~9, and the relative exceptional estimate (4.24), p.~11.
\end{corollary}
\begin{proof}
Choose $\theta=3/5$ and a fixed $C\ge C_Y(3/5)$.
Then $h>x^{7/12}$, $y<h<x$, and $y=o(h)$, exactly as required in the
first lemma. Its proposed formula would give
\[
 \frac1{M_x}\sum_{n\in S_x}\omega(n)=O(\loglog y).
\]
However,
\[
 \frac{u}{\loglog y}\asymp
 \frac{(\log x)^{1/3}}{(\loglog x)^{7/3}}\longrightarrow\infty,
\]
contradicting \eqref{eq:mean}. Changing the left endpoint of the interval
from open to closed adds at most one integer; its contribution is at most
$\log_2(2x)$ and is negligible relative to $M_x$.
The proposed second moment is $O(M_x(\loglog y)^2)$, whereas
\eqref{eq:tail} gives a lower bound
$\kappa^2u^2M_x(1-o(1))$ for any fixed $0<\kappa<3/5$.
Finally the threshold in (4.24) is bounded by a fixed power of
$\loglog x$. Theorem~\ref{thm:tail} shows that the proportion above it
tends to one, although the claimed relative upper bound tends to zero.
\end{proof}

\begin{remark}
This refutes those intermediate assertions, not $abc$ and not the possibility
of an exceptionally sparse counterexample family. In particular a set of
size $x^{o(1)}$ can still be infinite. It is also important to distinguish
$o(h)$ from $o(M_x)$: the whole smooth set already has size $o(h)$.
\end{remark}

\subsection{Two further transfer tests}
For a fixed integer $k\ge4$, the intervals
\[
 (p^k,p^k+(p^k)^{3/5}],\qquad X<p^k\le2X,
\]
are pairwise disjoint for large $X$. Indeed, successive distinct integer
bases give gaps at least $X^{1-1/k}$, larger than $(2X)^{3/5}$.
The finite encoding, using the largest thresholds in this range, shows that
at most $X^{o(1)}$ centres can contain a smooth candidate with
$\omega(n)\le2\loglog(p^k)$ at the scale \eqref{eq:parameters}.
The prime number theorem gives $X^{1/k+o(1)}$ centres, so their proportion is
at most $X^{-1/k+o(1)}$. This rules out a positive-proportion selection
of this particular kind, while leaving a sparse selection unresolved.

Recent work on exceptional sets \cite{BBLT,Bernert} and on the existence
of smooth neighbours \cite{Jain} addresses different quantifiers.
A useful general check is the following. If a counting theorem gives
$N(T)\ll T^{\eta+o(1)}$, then a seed of height $H$ contradicts it only
when it produces more than $H^{\beta\eta+o(1)}$ \emph{distinct} qualifying
outputs of height at most $O(H^\beta)$. Repeated encodings of one output do
not count. None of the existence results just cited supplies the required
uniform low-radical neighbours or such an amplification theorem.

\section{The coupled signed prime-exponent defect}\label{sec:signed}

For $n\ge1$ put
\[
 E_1(n)=\sum_{\substack{p\mid n\\v_p(n)=1}}\log p,
 \qquad
 E_3(n)=\sum_{p\mid n}\pospart{v_p(n)-2}\log p.
\]
Write $m=\min(a,b)$ and $M=\max(a,b)$ for an abc triple.

\begin{proposition}[Exact signed identity]\label{prop:signed}
For every positive integer $n$,
\begin{equation}\label{eq:single}
 \delta_2(n):=\log n-2\log\rad(n)=E_3(n)-E_1(n).
\end{equation}
The coupled defect satisfies
\begin{align}
 \Delta&:=\log(Mc)-2R\notag\\
 &=E_3(M)+E_3(c)-E_1(M)-E_1(c)-2\log\rad(m),\label{eq:coupled}\\
 2h-2R-\log2&\le\Delta\le2h-2R.\label{eq:corridor}
\end{align}
Consequently $abc$ is equivalent to
\begin{equation}\label{eq:coupledbound}
 \begin{split}
 \forall\eps>0\quad\exists K_\eps\quad\forall(a,b,c):\qquad
 E_3(M)+E_3(c)\le{}&E_1(M)+E_1(c)\\
 &+2\log\rad(m)+2\eps R+K_\eps.
 \end{split}
\end{equation}
\end{proposition}
\begin{proof}
For every positive exponent $e$,
$e-2=\pospart{e-2}-\mathbf1_{e=1}$. Multiply by $\log p$ and sum over the
actual prime factorization to obtain \eqref{eq:single}.
Pairwise coprimality gives
$r=\rad(m)\rad(M)\rad(c)$, which proves \eqref{eq:coupled}.
Since $c/2\le M\le c$, taking logarithms of
$c^2/2\le Mc\le c^2$ gives \eqref{eq:corridor}.
An $abc$ constant $C_\eps$ supplies $K_\eps=2C_\eps$.
Conversely \eqref{eq:coupledbound} and the lower corridor give
$C_\eps=(K_\eps+\log2)/2$.
\end{proof}

The equivalence does not prove \eqref{eq:coupledbound}; it identifies the
entire missing estimate without losing the negative exponent-one terms.
Replacing the coupled inequality by separate endpoint bounds fails.

\begin{proposition}[An unbounded separate-endpoint obstruction]\label{prop:dyadic}
For every $0\le\rho<1$ and $K\in\R$, there is an abc triple with $m=1$ and
\[
 \delta_2(M)>\log\rad(m)+\rho R+K.
\]
\end{proposition}
\begin{proof}
Take $(a,b,c)=(1,2^N,2^N+1)$ with $N\ge1$.
Then $m=1$, $M=2^N$, and $\delta_2(M)=(N-2)\log2$.
Without making any assumption about the odd neighbour,
\[
 r\le2(2^N+1)\le2^{N+2},\qquad R\le(N+2)\log2.
\]
Therefore
\[
 \delta_2(M)-\rho R\ge
 ((1-\rho)N-2-2\rho)\log2\longrightarrow\infty.
\]
This proves the claim. The bound uses only $\rad(n)\le n$ and
$\rad(2^N)=2$.
\end{proof}

This argument neither proves nor assumes that all odd neighbours have large
radical. It only rejects a false stronger estimate that discards their
possible contribution. The coupled route remains open.

\section{Finite certificates for reachable local lattices}\label{sec:lattice}

Let $K$ be a nonarchimedean local field, $\OO$ its valuation ring,
$\pi$ a uniformizer, and $k=\OO/\pi\OO$ its residue field, with
$|k|=q>1$ and $v(\pi)=1$. Fix $N\ge1$, write $L=\OO^N$, and normalize
additive Haar measure on $K^N$ by $\mu(L)=1$.
For any set $\mathcal R\subset L$, let
\[
 H(\mathcal R)=\overline{\Span_{\OO}(\mathcal R)}.
\]
If nonsingular $N$-tuples exist in $\mathcal R$, define
\[
 d(\mathcal R)=\min_{\substack{r_1,\ldots,r_N\in\mathcal R\\
                         \det[r_1\ \cdots\ r_N]\ne0}}
                         v(\det[r_1\ \cdots\ r_N]);
\]
otherwise set $d(\mathcal R)=\infty$.

\begin{theorem}[A finite witness for the exact hull measure]\label{thm:lattice}
With $q^{-\infty}=0$, one has
\[
 \mu(H(\mathcal R))=q^{-d(\mathcal R)}.
\]
If $d(\mathcal R)<\infty$, some matrix $A$ whose columns belong to
$\mathcal R$ satisfies
\[
 v(\det A)=d(\mathcal R),\qquad H(\mathcal R)=A\OO^N.
\]
\end{theorem}
\begin{proof}
In the full-rank case the nonzero determinant valuations are nonnegative
integers, so choose $A$ attaining their minimum $d$.
For $r\in\mathcal R$, Cramer's rule gives
\[
 (A^{-1}r)_i=\frac{\det A_i(r)}{\det A},
\]
where $A_i(r)$ replaces column $i$ by $r$.
Its columns still belong to $\mathcal R$. Every nonzero numerator has
valuation at least $d$; a zero numerator causes no difficulty.
Thus $A^{-1}r\in\OO^N$ and $\mathcal R\subset A\OO^N$.
Since the columns of $A$ lie in $\mathcal R$, their integral spans are equal.
The set $A\OO^N$ is compact, hence closed.

Integral row and column operations give
$UAV=\operatorname{diag}(\pi^{a_1},\ldots,\pi^{a_N})$, with
$U,V\in\GL_N(\OO)$, $a_i\ge0$, and $\sum a_i=d$.
One may obtain this diagonalization by repeatedly moving an entry of least
valuation to the upper left and clearing its row and column by integral
operations. The quotient $\OO^N/A\OO^N$ has $q^d$ elements.
Its equal-measure cosets partition $L$, proving $\mu(A\OO^N)=q^{-d}$.

In the rank-deficient case a nonzero $K$-linear functional annihilates
$\mathcal R$. Scale its coefficients to lie in $\OO$ with at least one
unit coefficient. It then defines a surjection $\ell:L\to\OO$ whose
closed kernel contains $H(\mathcal R)$. For every $t\ge1$,
$\ell^{-1}(\pi^t\OO)$ has index $q^t$ in $L$. Hence
$0\le\mu(H(\mathcal R))\le q^{-t}$ for all $t$, proving measure zero.
\end{proof}

\begin{corollary}\label{cor:residue}
The equality $H(\mathcal R)=L$ holds if and only if the reductions of
$\mathcal R$ span $k^N$, equivalently if some reachable column matrix has
unit determinant. More generally, any measurable $\OO$-submodule containing
all columns of a reachable matrix $A$ with $v(\det A)=d$ has measure at least
$q^{-d}$.
\end{corollary}
\begin{proof}
A determinant is a unit exactly when its reduction is nonzero; a spanning
set of $k^N$ contains a basis. Apply Theorem~\ref{thm:lattice}.
For the last claim, the target contains $A\OO^N$ and measure is monotone.
The target need not be contained in $L$.
\end{proof}

\subsection{Independent and simultaneous actions}
Let $L_j=\OO^{d_j}$, with $d_j\ge1$, and let
$T=\bigotimes_{j=1}^s L_j$, where $s\ge1$ and
$N=\prod_jd_j$. For nonzero $x_j\in L_j$, write
$x_j=\pi^{m_j}u_j$ with $u_j$ primitive.

\begin{proposition}\label{prop:independent}
Under the independent action of $\prod_j\GL_{d_j}(\OO)$,
the orbit of $\bigotimes_jx_j$ spans
$\pi^{\sum_jm_j}T$. The measure of this spanned lattice, normalized by
$\mu(T)=1$, is
$q^{-N\sum_jm_j}$.
\end{proposition}
\begin{proof}
Integral row operations send a primitive vector to any chosen standard
basis vector. Choosing these operations independently in each factor gives
every standard tensor basis vector multiplied by $\pi^{\sum_jm_j}$.
Conversely all orbit vectors lie in this scaled lattice.
Its index is $q^{N\sum_jm_j}$, giving the measure.
\end{proof}

\begin{proposition}[A strict simultaneous-action counterexample]\label{prop:diagonal}
Over any commutative ring $A$, identify $A^2\otimes_AA^2$ with
$\Mat_2(A)$. The $A$-span of $\{v\otimes v:v\in A^2\}$ is exactly
the submodule of symmetric matrices. If $A\ne0$, it omits $E_{12}$.
For $A=\OO$ it has Haar measure zero in $\Mat_2(K)$.
\end{proposition}
\begin{proof}
Every generator has equal off-diagonal entries, a condition preserved by
linear combinations. Conversely $e_1\otimes e_1$, $e_2\otimes e_2$, and
\[
 (e_1+e_2)\otimes(e_1+e_2)-e_1\otimes e_1-e_2\otimes e_2
 =E_{12}+E_{21}
\]
span all symmetric matrices. No division by two occurs.
The functional $M\mapsto M_{12}-M_{21}$ vanishes on the span and takes
value one on $E_{12}$. In the local-field case, the congruence
$M_{12}=M_{21}\pmod{\pi^t}$ occupies a proportion $q^{-t}$ of the
integral matrix lattice for every $t$. Its exact equality locus therefore
has measure zero.
\end{proof}

The generating set in Proposition~\ref{prop:diagonal} is stable under every
simultaneous action $g\otimes g$. Thus even the full simultaneous integral
linear group does not provide Proposition~\ref{prop:independent}.
This rejects a specific abstract substitute for independent generation; it
does not identify the actual IUT outputs with this diagonal example.

\subsection{What a source-derived application must establish}
The July 2026 LANA report \cite{LANA} isolates an identification between
q-pilot and possible-output volume data in the passage to IUT~III,
Corollary~3.12 \cite{IUT3}. Joshi's July 30 response \cite{JoshiResponse}
refers to distinct arithmeticoids and his Part~IV comparison
\cite[Theorem~6.10.1]{JoshiIV}. These are source claims to reconstruct,
not assumptions used in our local theorems.

Part~III \cite[Proposition~9.7.5.1]{JoshiIII} allows specified topological
isomorphisms in its collation. It would therefore be inaccurate to say
that it permits no automorphism enlargement. However, the simultaneous
actions and the product-to-tensor hull operation in
Theorem-Definition~9.8.1.1 do not, by themselves, state unrestricted
independent integral linear actions. One must prove which choices are
permitted and that the actual target contains their integral span.

For a source-faithful finite packet of local sets
$\mathcal R_v\subset\OO_v^{N_v}$ and positive weights $\lambda_v$,
Theorem~\ref{thm:lattice} gives, when all defects are finite,
\[
 V_H=-\sum_v\frac{\lambda_v}{N_v}d(\mathcal R_v)\log q_v.
\]
Suppose genuine local target regions $\Theta_v$ contain the respective
hulls and have finite positive measure. Define, on these same carriers,
$V_\Theta=\sum_v(\lambda_v/N_v)\log\mu_v(\Theta_v)$.
Let $V_q\in\R$ be the fixed native $q$-pilot log-volume under the same
source normalization. Then
\[
 \sum_v\frac{\lambda_v}{N_v}d(\mathcal R_v)\log q_v\le -V_q
 \quad\Longrightarrow\quad V_q\le V_H\le V_\Theta.
\]
The final inequality follows from local measure monotonicity, followed by
the logarithm and the positive weighted sum. A different volume functional
on a non-product region requires a separate monotonicity theorem.
This is a sufficient numerical interface. Constructing the genuine vectors,
establishing closure and normalization, controlling the determinant sum,
and deriving the uniform IUT~IV height comparison \cite{IUT4} remain
separate obligations. Archimedean factors and infinite-packet convergence
are not covered by the finite local theorem.

\section{Effective concentration in a specified Pell residual}\label{sec:pell}

We now restrict to positive integers $b,A,B,u,v,r_0,s,X$ and a prime
$p\ge37$ satisfying
\begin{align}
 b=Au^2,\quad b+1=Bv^2,\quad b+2=3r_0^2,\quad b+3=s^2,\label{eq:packet1}\\
 D=3AB,\quad Z=b^2+3b+1=T_p(X),\quad
 D+1\le X^2,\quad X^p\le Z.\label{eq:packet2}
\end{align}
Here $T_p$ is the Chebyshev polynomial of the first kind. These equations
define the specified residual packet; no reduction of general abc
counterexamples to it is asserted.
Set
\[
 H=\log(b+2),\quad L=\log(D+1),\quad
 C_0=12^{331776},\quad K_0=\log(2C_0).
\]
The packet immediately yields
\begin{equation}\label{eq:packetgate}
 (D+1)^p\le X^{2p}\le Z^2<(b+2)^4,\qquad pL<4H.
\end{equation}

Theorem~2.2 of B\'erczes--Evertse--Gy\H{o}ry \cite{BEG}, specialized to
$K=\Q$, $S=\{\infty\}$, and exponent two, gives the following published
input. If $f\in\Z[t]$ is squarefree of degree $n\ge3$, with coefficient
height $\widehat h=\log\max(1,|a_0|,\ldots,|a_n|)$, and $f(x)=y^2$
in integers, then
\begin{equation}\label{eq:beg}
 \max\{\height(x),\height(y)\}
 \le (4n)^{4096n^4}\exp(50n^4\widehat h),
 \quad \height(t)=\log\max(1,|t|).
\end{equation}

\begin{proposition}[Two effective constraints]\label{prop:pell}
Every packet \eqref{eq:packet1}--\eqref{eq:packet2} satisfies
\begin{equation}\label{eq:two-effective}
 H\le\exp(4300p^5),\qquad H\le2C_0D^{4050}.
\end{equation}
\end{proposition}
\begin{proof}
First put $f_p(t)=4T_p(t)+5$. The relation in \eqref{eq:packet2} gives
$f_p(X)=(2b+3)^2$.
The identities $T'_p=pU_{p-1}$ and
$T_p^2-1=(t^2-1)U_{p-1}^2$ show that a repeated root of $f_p$ would have
$T_p=\pm1$ and $T_p=-5/4$, a contradiction.
If $M_n$ is the sum of the absolute values of the coefficients of $T_n$,
then $M_{n+2}\le2M_{n+1}+M_n$ and $M_0=M_1=1$.
Induction gives $M_n\le3^n$. Thus every coefficient of $f_p$ has
absolute value at most $4\cdot3^p+5\le3^{p+2}$, and
$\widehat h\le(p+2)\log3\le4p$.
Since $\log(4p)\le p$ for $p\ge3$, \eqref{eq:beg} implies
\[
 H\le\log(2b+3)
 \le\exp\!\left(4096p^4\log(4p)+50p^4\widehat h\right)
 \le\exp(4300p^5).
\]

For the second estimate, set $x=b+1$ and $Y=Duvr_0$.
The first three equations in \eqref{eq:packet1} give
\begin{equation}\label{eq:ellipticpacket}
 Y^2=D(x^3-x).
\end{equation}
The cubic $D(t^3-t)$ is squarefree and has coefficient height $\log D$.
Equation~\eqref{eq:beg}, with $n=3$, gives
\[
 \log(b+1)\le12^{4096\cdot3^4}\exp(50\cdot3^4\log D)
 =C_0D^{4050}.
\]
Finally $b+2\le(b+1)^2$ for $b\ge1$, proving the second bound.
\end{proof}

\begin{proposition}[Descent to the squarefree parameter]\label{prop:descent}
Write $D=d\tau^2$, where $d$ is a positive squarefree integer and
$\tau\ge1$ is an integer. Then every packet satisfies
\[
 H\le2C_0d^{4050}.
\]
\end{proposition}
\begin{proof}
Equation~\eqref{eq:ellipticpacket} implies $\tau^2\mid Y^2$.
Comparing prime valuations gives $\tau\mid Y$. Write $Y=\tau y$ and
cancel the nonzero factor $\tau^2$ to obtain
\[
 y^2=d(x^3-x),\qquad x=b+1.
\]
The polynomial $d(t^3-t)$ has distinct roots and coefficient height
$\log d$. Equation~\eqref{eq:beg} gives
$\log(b+1)\le C_0d^{4050}$, and again $H\le2\log(b+1)$.
\end{proof}

Combining \eqref{eq:packetgate} and \eqref{eq:two-effective} gives
\begin{align}
 p&\ge\left(\frac{\log H}{4300}\right)^{1/5},\qquad
 L<4\cdot4300^{1/5}\frac{H}{(\log H)^{1/5}},\label{eq:uppercorr}\\
 D&\ge\left(\frac{H}{2C_0}\right)^{1/4050},\qquad
 p(\log H-K_0)<16200H.\label{eq:lowercorr}
\end{align}
For example, the first upper bound follows by multiplying the lower bound
for $p$ by $L>0$ and applying $pL<4H$.
For the final inequality use
$\log H\le K_0+4050\log D$ and $\log D<L$.
In a form avoiding fractional powers, the former calculation is
\[
 (\log H)L^5\le4300(4H)^5.
\]

\begin{corollary}\label{cor:residualsequence}
Any sequence of these packets with strictly increasing $b_n$, with $d_n$
the positive squarefree representative of $D_n$, satisfies
\[
 p_n\longrightarrow\infty,\quad d_n\longrightarrow\infty,
 \quad D_n\longrightarrow\infty,
 \quad\frac{\log(D_n+1)}{\log(b_n+2)}\longrightarrow0.
\]
\end{corollary}
\begin{proof}
Here $H_n\to\infty$ and $\log H_n\to\infty$.
Apply \eqref{eq:uppercorr} and \eqref{eq:lowercorr}.
The bound in Proposition~\ref{prop:descent} forces $d_n\to\infty$.
\end{proof}

These restrictions are compatible, so they do not exclude the residual
sequence. Distinct positive squarefree integers $d>1$ give distinct
quadratic fields $\Q(\sqrt d)$, so these fields cannot remain in a fixed
finite collection.
Also $\log d_n/H_n\to0$ because $1\le d_n\le D_n$.
Thus freezing either the index or the quadratic field cannot be justified
for an unbounded sequence. The full argument
uses \eqref{eq:beg}; Lean checks only the explicit arithmetic and scalar
consequences with that input stated as a hypothesis.

\section{Applicability of recent elliptic-curve estimates}\label{sec:frey}

\subsection{A denominator lower bound}
For the Frey curve $E:y^2=x(x-a)(x+b)$ attached to an abc triple, let
$Q=a^2+ab+b^2$. Its invariant is
\[
 j_E=\frac{256Q^3}{(abc)^2}.
\]
Let $n,d$ be the positive numerator and denominator in lowest terms.

\begin{proposition}\label{prop:denominator}
For every such triple,
\begin{equation}\label{eq:denominator}
 c^4\le1024d,\qquad n\le256c^6.
\end{equation}
\end{proposition}
\begin{proof}
Primitivity implies $\gcd(Q,abc)=1$. Indeed, reduction modulo a prime
dividing any one of $a,b,c$ leaves in $Q$ the square of a nonzero other
coordinate. Therefore $g=\gcd(256Q^3,(abc)^2)$ divides $256$, and
$(abc)^2=gd\le256d$.
Since $(a-1)(b-1)\ge0$, we have $ab\ge c-1$; as $c\ge2$, this gives
$2ab\ge c$. Squaring $2abc\ge c^2$ proves
$c^4\le4(abc)^2\le1024d$.
Finally $Q=c^2-ab\le c^2$, so $n\le256Q^3\le256c^6$.
\end{proof}

Pasten's 2026 small-denominator theorem
\cite[Theorem~1.2]{PastenSmall} applies under a bound
$\den(j_E)\le A_0(\log\num(j_E))^{B_0}$ with fixed positive
$A_0,B_0$. For Frey curves, \eqref{eq:denominator} would imply
\[
 c^4\le1024A_0(\log256+6\log c)^{B_0}.
\]
A positive power of $c$ dominates every fixed power of its logarithm.
Thus only bounded $c$, and consequently finitely many positive triples,
can satisfy this hypothesis for fixed $A_0,B_0$.
This is a precise restriction on direct applicability, not a
counterexample to the theorem or to other Szpiro strategies.

\subsection{The August 2026 Mordell estimate}
The preprint \cite[Theorem~1.3]{PastenNew}, submitted August 24, 2026,
gives an effective bound for $y^2=x^3+k$, $k\ne0$, using
\[
 \underline{k}=\prod_{q\mid k}q^{\min(2,v_q(k))}.
\]
The stated estimate is
\begin{equation}\label{eq:pasten}
 \log\max(|x|,|y|)\ll
 \underline{k}^{1/2}\log^2(2\underline{k})\log(2|k|)
 \log\!\left(\underline{k}\log(3|k|)\right).
\end{equation}
We record this new result without treating its publication date as a
substitute for checking its hypotheses or proof.

Directly substituting $c_6^2=c_4^3-1728\Delta$ and
$\Delta=16(abc)^2$ gives $k=-27648(abc)^2$.
Then $\underline{k}^{1/2}\ll r$, while both
$\log(2|k|)$ and $\log\max(|c_4|,|c_6|)$ are comparable with $h$.
The resulting right-hand side has the shape
\[
 r\log^2(2r)\,h\,\log(2r\log(3c)).
\]
The height remains a multiplicative factor on the right; cancellation
does not yield an upper bound for $c$ in terms of $r$.
The integral-$j$ result in Theorem~1.6 of the same paper concerns the
height of the \emph{curve}, and cannot be substituted for a point-height bound.

There is also a strict obstruction to one possible transfer from
\eqref{eq:ellipticpacket}. With $U=Dx$ and $V=DY$, it becomes
$E_D:V^2=U^3-D^2U$, of $j$-invariant $1728$.

\begin{proposition}\label{prop:cm}
For $D,k\ne0$, there is no nonconstant algebraic morphism over
$\overline{\Q}$ from $E_D$ to the nonsingular Mordell curve
$M_k:v^2=u^3+k$.
\end{proposition}
\begin{proof}
Their geometric endomorphism algebras are respectively $\Q(i)$ and
$\Q(\sqrt{-3})$: the automorphisms
$(x,y)\mapsto(-x,iy)$ and $(u,v)\mapsto(\zeta_3u,v)$ give these
quadratic actions, and in characteristic zero the endomorphism algebra
of an elliptic curve has dimension at most two over $\Q$.
A nonconstant morphism, translated to send the origin to the origin,
is an isogeny. An isogeny induces an isomorphism of rational endomorphism
algebras by conjugation. The two imaginary quadratic fields are not
isomorphic, a contradiction; see also \cite[Chapters III and VI]{Silverman}.
\end{proof}

The proposition excludes this morphism strategy, even if $k$ depends on
$D$. It does not exclude other Diophantine encodings or extending the new
logarithmic-form methods to a different family.

\section{Formal verification and remaining obligations}\label{sec:formal}

The companion development uses Lean~4.32.0 and the repository's pinned
dependencies. The new files are:
\begin{center}
\begin{tabular}{>{\raggedright\arraybackslash}p{0.45\textwidth}>{\raggedright\arraybackslash}p{0.45\textwidth}}
\toprule
Module (within \texttt{IUTThreeClosures}) & Checked mathematical scope\\
\midrule
\path{SmoothLowOmegaCounting} & Actual prime-power encoding, finite
cardinality bounds, and the finite first moment.\\[3pt]
\path{SignedPrimeSupport} & Actual prime-exponent sums and the coupled
equivalence with the original target.\\[3pt]
\path{SignedEndpointDyadicObstruction} & The genuine dyadic triples and
failure of a separate-endpoint slope below one.\\[3pt]
\path{ReachableTensorLatticeCriterion} & Unit-determinant span certificate,
symmetric span and its missing off-diagonal tensor.\\[3pt]
\texttt{FreyPellEffectiveResidual\allowbreak Corridor20260830} & The integral elliptic
identity, integer square-factor descent, explicit scalar implications,
and the Frey denominator bound.\\
\bottomrule
\end{tabular}
\end{center}

These modules introduce no new mathematical axiom. The validation record
lists their exact commands, results, and kernel dependency reports.
The pre-existing signed-layer module required an explicit import, closure
of an elementary algebra branch, and the correct order of unfolding before
rewriting; these were repaired without changing its statements.

The formal scope is deliberately narrower than the full paper.
Theorem~\ref{thm:tail} and Corollary~\ref{cor:carella} use Younis and classical
smooth-number asymptotics in their mathematical proofs; those external
analytic theorems are not formalized here. The complete local Haar formula
in Theorem~\ref{thm:lattice}, the general independent-action theorem,
the effective bound \eqref{eq:beg}, the asymptotic consequences of the
Pell bounds, and the endomorphism-algebra argument also retain their
stated paper-proof boundary. A Lean implication with an explicit height
hypothesis is not a Lean proof of that hypothesis.

The surviving proof target is a uniform estimate for the coupled expression
\eqref{eq:coupledbound}, or a source-derived theorem strong enough to imply
it. The IUT route still needs genuine reachable output data, compatible
normalizations and the global comparison. The analytic disproof route still
needs a height-unbounded family satisfying one fixed radical exponent gap;
the finite-support bound neither constructs nor forbids such a sparse
family. The Pell residual inequalities currently admit a nonempty numerical
corridor and do not apply to all abc triples. No unconditional closed term
of type \texttt{ABCConjecture}, or of its negation, is produced by this work.

\begin{thebibliography}{99}

\bibitem{BEG}
A.~B\'erczes, J.-H.~Evertse and K.~Gy\H{o}ry,
Effective results for hyper- and superelliptic equations over number fields,
\emph{Publ. Math. Debrecen} \textbf{82} (2013), 727--756.
\href{https://doi.org/10.5486/PMD.2013.5748}{doi:10.5486/PMD.2013.5748}.

\bibitem{BBLT}
C.~Bernert, T.~Browning, J.~D.~Lichtman and J.~Ter\"av\"ainen,
Bounds on the exceptional set in the $abc$ conjecture,
\href{https://arxiv.org/abs/2410.12234v2}{arXiv:2410.12234v2}.

\bibitem{Bernert}
C.~Bernert, The exceptional set in the $abc$ conjecture,
\href{https://arxiv.org/abs/2506.13364v1}{arXiv:2506.13364v1}, 2025.

\bibitem{Carella}
N.~A.~Carella, Note on the Exceptional Set in the ABC Conjecture,
\href{https://arxiv.org/abs/2608.16764v2}{arXiv:2608.16764v2},
revised August 24, 2026.

\bibitem{HT}
A.~Hildebrand and G.~Tenenbaum, Integers without large prime factors,
\emph{J. Th\'eor. Nombres Bordeaux} \textbf{5} (1993), 411--484.
\href{https://doi.org/10.5802/jtnb.101}{doi:10.5802/jtnb.101}.

\bibitem{Jain}
S.~Jain, Existence of Smooth Numbers in Short Intervals,
\emph{Q. J. Math.} \textbf{77} (2026), no.~2, 397--422.
\href{https://doi.org/10.1093/qmath/haag010}{doi:10.1093/qmath/haag010}.

\bibitem{JoshiIII}
K.~Joshi, Construction of Arithmetic Teichm\"uller Spaces III:
A Rosetta Stone and a proof of Mochizuki's Corollary 3.12,
\href{https://arxiv.org/abs/2401.13508v4}{arXiv:2401.13508v4}, 2025.

\bibitem{JoshiIV}
K.~Joshi, Construction of Arithmetic Teichm\"uller Spaces IV:
Proof of the abc-conjecture,
\href{https://arxiv.org/abs/2403.10430v2}{arXiv:2403.10430v2}, 2025.

\bibitem{JoshiResponse}
K.~Joshi, Comments on the LANA Project Report of Kato et al.,
July 30, 2026, six pages.
\href{https://sites.arizona.edu/kirti-joshi/files/2026/07/Comments-on-the-LANA-Project-Report-of-Kato-et-al.pdf}{Author's original PDF}.

\bibitem{LANA}
F.~Kato et al., LANA report, revision 1.0.1,
\href{https://github.com/katobungen/LANA_report_202607/tree/293bdd89463473ae13d40834d70fb4b7ba81da1f}{pinned source revision}, July 20, 2026.
See also the \href{https://zen.ac.jp/news/zmcpostevent0717e}{ZEN University
interim report of July 17, 2026}.

\bibitem{IUT3}
S.~Mochizuki, Inter-universal Teichm\"uller Theory III:
Canonical Splittings of the Log-Theta-Lattice,
\emph{Publ. Res. Inst. Math. Sci.} \textbf{57} (2021), 403--626.
\href{https://doi.org/10.4171/PRIMS/57-1-3}{doi:10.4171/PRIMS/57-1-3}.

\bibitem{IUT4}
S.~Mochizuki, Inter-universal Teichm\"uller Theory IV:
Log-volume Computations and Set-theoretic Foundations,
\emph{Publ. Res. Inst. Math. Sci.} \textbf{57} (2021), 627--723.
\href{https://doi.org/10.4171/PRIMS/57-1-4}{doi:10.4171/PRIMS/57-1-4}.

\bibitem{PastenNew}
H.~Pasten, Power-saving bounds for Thue--Mahler and Mordell equations,
\href{https://arxiv.org/abs/2608.23559v1}{arXiv:2608.23559v1}, August 24, 2026.

\bibitem{PastenSmall}
H.~Pasten, Szpiro's conjecture when the denominator of the $j$-invariant
is small, \emph{Cubo} \textbf{28} (2026), no.~2, 383--389.
\href{https://doi.org/10.56754/0719-0646.2802.383}{doi:10.56754/0719-0646.2802.383}.

\bibitem{Silverman}
J.~H.~Silverman, \emph{The Arithmetic of Elliptic Curves},
2nd ed., Graduate Texts in Mathematics 106, Springer, 2009.

\bibitem{Younis}
K.~Younis, Asymptotics for smooth numbers in short intervals,
\href{https://arxiv.org/abs/2409.05761v1}{arXiv:2409.05761v1}, 2024.

\end{thebibliography}
\end{document}
